%%--------------------------------------------------------------------
%1--------------------------------------------------------------------
\begin{glyph}{Floor (階)}{floor}\label{glyph:floor}
Floor, story (of a building).\\\hline
Related to this idea of layers, 階 can also refer to \textit{rank} or \textit{social stratum} to which a person or group of people belong to.
\end{glyph}
\gindex{Floor}{階}\strindex{12}{階}

\begin{comps}
\comprtwocone&阝阜 + 皆\\\hline
\comprthreecone&Hill/Mound + All (p. \pageref{glyph:all})\\\hline
\end{comps}

\begin{remglyph}
\hooglig{The level} of the \remcom{hill} where \remcom{everyone} is.\\\hline
\end{remglyph}

\begin{prons}
\pronsrtwocone&\textbf{カイ (KAI)}&---\\\hline
\pronsrthreecone&---&---\\\hline
\end{prons}

\begin{remprons}
\hooglig{Floor} selling \comon{kites}.\\\hline
\end{remprons}

%{\middel}
% &  \mbox{()}\\\hline
\begin{somme}
一階&one + floor = \textit{first floor; ground floor}&イッ{\middel}カイ \mbox{(IK{\middel}KAI)}\\\hline
階上&floor + upper = \textit{upstairs}&カイ{\middel}ジョウ \mbox{(KAI{\middel}J\={O})}\\\hline 
階下&floor + lower = \textit{downstairs}&カイ{\middel}カ \mbox{(KAI{\middel}KA)}\\\hline
下の階&lower + floor = \textit{lower floor}&した{\middel}の{\middel}カイ \mbox{(shita{\middel}no{\middel}KAI)}\\\hline
最上階&best + floor = \textit{top floor}&サイ{\middel}ジョウ{\middel}カイ \mbox{(SAI{\middel}J\={O}{\middel}KAI)}\\\hline
地階&ground + floor = \textit{basement; cellar}&チ{\middel}カイ \mbox{(CHI{\middel}KAI)}\\\hline
二階建て&2 + floor + construct = \textit{two-storied building}&ニ{\middel}カイ{\middel}だ{\middel}て \mbox{(NI{\middel}KAI{\middel}da{\middel}te)}\\\hline
二階屋&upstairs + roof = \textit{two-storey house}&ニ{\middel}カイ{\middel}や \mbox{(NI{\middel}KAI{\middel}ya)}\\\hline
音階&sound + floor = \textit{music scale}&オン{\middel}カイ \mbox{(ON{\middel}KAI)}\\\hline
五音音階&pentatonic + scale = \textit{pentatonic scale}&ゴ{\middel}オン{\middel}オン{\middel}カイ \mbox{(GO{\middel}ON{\middel}ON{\middel}KAI)}\\\hline
長音階&long sound + floor = \textit{major scale}&チョウ{\middel}オン{\middel}カイ \mbox{(CH\={O}{\middel}ON{\middel}KAI)}\\\hline
短音階&short sound + floor = \textit{minor scale}&タン{\middel}オン{\middel}カイ \mbox{(TAN{\middel}ON{\middel}KAI)}\\\hline
半音階&half-tone + floor = \textit{chromatic scale}&ハン{\middel}オン{\middel}カイ \mbox{(HAN{\middel}ON{\middel}KAI)}\\\hline
全音階&whole tone + floor = \textit{diatonic scale}&ゼン{\middel}オン{\middel}カイ \mbox{(ZEN{\middel}ON{\middel}KAI)}\\\hline
各階&each + floor = \textit{each floor}&カク{\middel}カイ \mbox{(KAKU{\middel}KAI)}\\\hline
表二階&express + upstairs = \textit{second-floor front room}&おもて{\middel}ニ{\middel}カイ \mbox{(omote{\middel}NI{\middel}KAI)}\\\hline
\end{somme}
\end{multicols}


\begin{decomp}{7}
その&かばん&を&階上&に&運び&なさい。\\\hline
その&かばん&を&カイジョウ&に&はこび&なさい\\\hline
that/the&bag&を&upstairs&に&to carry&to do\\\hline
\multicolumn{7}{|r|}{\textit{Carry the bags upstairs.}}\\\hline
\end{decomp}
\pagebreak
%%--------------------------------------------------------------------
%2--------------------------------------------------------------------
\begin{multicols}{2}
\begin{glyph}{Point (点)}{point}\label{glyph:point}
Point.  Mark.
\end{glyph}
\gindex{Point}{点}\strindex{9}{点}

\begin{comps}
\comprtwocone&占 + 火灬\\\hline
\comprthreecone&Occupy (p. \pageref{glyph:occupy}) + Fire\\\hline
\end{comps}

\begin{remglyph}
\hooglig{Points} of \remcom{fortune telling} under \remcom{fire}.\\\hline
\end{remglyph}

\begin{prons}
\pronsrtwocone&\textbf{テン (TEN)}&---\\\hline
\pronsrthreecone&---&---\\\hline
\end{prons}

\begin{remprons}
\hooglig{Points} cause a lot of \comon{tension}.\\\hline
\end{remprons}

%{\middel}
% &  \mbox{()}\\\hline
\begin{somme}
点火&point + fire = \textit{ignition}&テン{\middel}カ \mbox{(TEN{\middel}KA)}\\\hline
重点&heavy + point = \textit{important point; emphasis}&ジュウ{\middel}テン \mbox{(J\={U}{\middel}TEN)}\\\hline
氷点&ice + point = \textit{freezing point}&ヒョウ{\middel}テン \mbox{(HY\={O}{\middel}TEN)}\\\hline
欠点&lack + point = \textit{fault; defect}&ケッ{\middel}テン \mbox{(KET{\middel}TEN)}\\\hline
弱点&weak + point = \textit{weakness}&ジャク{\middel}テン \mbox{(JAKU{\middel}TEN)}\\\hline
要点&essential + point = \textit{gist; main point}&ヨウ{\middel}テン \mbox{(Y\={O}{\middel}TEN)}\\\hline
地点&ground + point = \textit{site; point on a map}&チ{\middel}テン \mbox{(CHI{\middel}TEN)}\\\hline
終点&finish + point = \textit{terminus; last stop}&シュウ{\middel}テン \mbox{(SH\={U}{\middel}TEN)}\\\hline
原点&original + point = \textit{origin; starting point}&ゲン{\middel}テン \mbox{(GEN{\middel}TEN)}\\\hline
\end{somme}
\end{multicols}

\begin{decomp}{6}
気温&は&氷点下&６&度&です。\\\hline
キオン&は&ヒョウテンカ&６&ド&です\\\hline
temperature&は&below freezing&6&degree&be/is\\\hline
\multicolumn{6}{|r|}{\textit{It's six degrees below zero.}}\\\hline
\end{decomp}

%%--------------------------------------------------------------------
%%--------------------------------------------------------------------
\begin{multicols}{2}
\begin{glyph}{Side (側)}{side}\label{glyph:side}
Side.
\end{glyph}
\gindex{Side}{側}\strindex{11}{側}

\begin{comps}
\comprtwocone& 人亻 + 則\\\hline
\comprthreecone&Man/Person + Rule (p. \pageref{glyph:rule})\\\hline
\end{comps}

\begin{remglyph}
\hooglig{Side} with \remcom{people} who \remcom{rule}.\\\hline
\end{remglyph}

\begin{prons}
\pronsrtwocone&\textbf{ソク (SOKU)}&\textbf{がわ (gawa)}\\\hline
\pronsrthreecone&---&---\\\hline
\end{prons}

\begin{remprons}
This \hooglig{side}, you'll find \comkun{guavas} stuffed in \comon{socks}.\\\hline
\end{remprons}

%{\middel}
% &  \mbox{()}\\\hline
\begin{somme}
側&\textit{side}&がわ \mbox{(gawa)}\\\hline
左側&left + side = \textit{left side}&ひだり{\middel}がわ \mbox{(hidari{\middel}gawa)}\\\hline
側近&side + near = \textit{close associate}&ソッ{\middel}キン \mbox{(SOK{\middel}KIN)}\\\hline
両側&both + side = \textit{both sides}&リョウ{\middel}がわ \mbox{(RY\={O}{\middel}gawa)}\\\hline
内側&inside + side = \textit{inside; inner part}&うち{\middel}がわ \mbox{(uchi{\middel}gawa)}\\\hline
外側&outside + side = \textit{exterior; outside}&そと{\middel}がわ \mbox{(soto{\middel}gawa)}\\\hline
受信側&reception + side = \textit{receiving entity}&ジュ{\middel}シン{\middel}がわ \mbox{(JU{\middel}SHIN{\middel}gawa)}\\\hline
発信側&despatch + side = \textit{sender; transmitter}&ハッ{\middel}シン{\middel}がわ \mbox{(HAS{\middel}SHIN{\middel}gawa)}\\\hline
不側不離&not + side + inseparability = \textit{close relationship}&フ{\middel}ソク{\middel}フ{\middel}リ \mbox{(FU{\middel}SOKU{\middel}FU{\middel}RI)}\\\hline
\notyet{反対側}&opposition + side = \textit{opposite side}&ハン{\middel}タイ{\middel}がわ \mbox{(HAN{\middel}TAI{\middel}gawa)}\\\hline
\end{somme}
\end{multicols}

\begin{decomp}{7}
日本&で&は&車&は&左側&です。\\\hline
ニッポン&で&は&くるま&は&ひだりがわ&です\\\hline
Japan&で&は&car&は&left side&be/is\\\hline
\multicolumn{7}{|r|}{\textit{In Japan, we drive on the left side of the road.}}\\\hline
\end{decomp}
\pagebreak
%%--------------------------------------------------------------------
%3--------------------------------------------------------------------

\begin{multicols}{2}
\begin{glyph}{Official (官)}{official}\label{glyph:official}
Official.  Government.\\\hline
As a suffix, this 漢字 indicates an \textit{official with governmental or bureaucratic authority}.
\end{glyph}
\gindex{Official}{官}\strindex{8}{官}

\begin{comps}
\comprtwocone&宀 + 㠯\\\hline
\comprthreecone&Roof + Two Mouths\\\hline
\end{comps}

\begin{remglyph}
\hooglig{The officials} under that \remcom{roof} are \remcom{two vampires} working for the \remcom{government}.\\\hline
\end{remglyph}

\begin{prons}
\pronsrtwocone&\textbf{カン (KAN)}&---\\\hline
\pronsrthreecone&---&---\\\hline
\end{prons}

\begin{remprons}
\hooglig{Officials} are very \comon{cunning}.\\\hline
\end{remprons}

%{\middel}
% &  \mbox{()}\\\hline
\begin{somme}
教官&teach + official = \textit{instructor}&キョウ{\middel}カン \mbox{(KY\={O}{\middel}KAN)}\\\hline
試験管&examination + official = \textit{examiner}&シ{\middel}ケン{\middel}カン \mbox{(SHI{\middel}KEN{\middel}KAN)}\\\hline
外交官&foreign policy + official = \textit{diplomat}&ガイ{\middel}コウ{\middel}カン \mbox{(GAI{\middel}K\={O}{\middel}KAN)}\\\hline
官話&official + speak = \textit{Mandarin}&カン{\middel}ワ \mbox{(KAN{\middel}WA)}\\\hline
高官&tall + official = \textit{high official}&コウ{\middel}カン \mbox{(K\={O}{\middel}KAN)}\\\hline
係官&offical + connection = \textit{official in charge}&かかり{\middel}カン \mbox{(kakari{\middel}KAN)}\\\hline
士官&gentleman + official = \textit{military officer}&シ{\middel}カン \mbox{(SHI{\middel}KAN)}\\\hline
書記官&secretary + official = \textit{secretary}&ショ{\middel}キ{\middel}カン \mbox{(SHO{\middel}KI{\middel}KAN)}\\\hline
官界&official + world = \textit{bureaucracy}&カン{\middel}カイ \mbox{(KAN{\middel}KAI)}\\\hline
\end{somme}
\end{multicols}

\begin{decomp}{8}
教官&は&私&に&毎日&運動する&様に&勧めた。\\\hline
キョウカン&は&わたし&に&マイニチ&ウンコウする&ヨウに&すすめた\\\hline
instructor&は&I&に&daily&exercise&ように&to recommend\\\hline
\multicolumn{8}{|r|}{\textit{The instructor advised me to ge exercise every day.}}\\\hline
\end{decomp}
%%--------------------------------------------------------------------
%4--------------------------------------------------------------------
\begin{multicols}{2}
\begin{glyph}{Feather (羽)}{feather}\label{glyph:feather}
Feather.\\\hline
This 漢字 also can also refer to the \textit{wings} of birds, insects and aeroplanes, etc.
\end{glyph}
\gindex{Feather}{羽}\strindex{6}{羽}

\begin{comps}
\comprtwocone&羽\\\hline
\comprthreecone&Feather\\\hline
\end{comps}

\begin{remglyph}
Part of the 漢字 radicals (\#{124}).\\\hline
\end{remglyph}

\begin{prons}
\pronsrtwocone&---&\textbf{は、はね (ha, hane)}\\\hline
\pronsrthreecone&ウ (U)&---\\\hline
\rye{3}{はね tends to be the reading used when the 漢字 appears on its own, while は is more common in compounds.}
\end{prons}

\begin{remprons}
\hooglig{Wings} come to a \comkun{halt} when exposed to \ucomon{ooky} \comkun{hanepoot}.\\\hline
{\warn}\textit{Hanepoot} is a variety of grape, grown in South Africa.\\\hline
\end{remprons}

%{\middel}
% &  \mbox{()}\\\hline
\begin{somme}
\rye{3}{The final compound uses the less common 音読み of 合 (カッ/KA-); this is the only time 合 is read with this pronunciation.}
羽&\textit{feather; wing}&はね \mbox{(hane)}\\\hline
{\diff}合羽&match + feather = \textit{raincoat}&カッ{\middel}ぱ \mbox{(KAP{\middel}pa)}\\\hline
羽毛&feather + hair/fur = \textit{feathers; plumage}&ウ{\middel}モウ \mbox{(U{\middel}M\={O})}\\\hline
天の羽衣&heaven + feather + clothes = \textit{angel's raiment}&は{\middel}ごろも \mbox{(ha{\middel}goromo)}\\\hline
手羽&hand + feather = \textit{chicken wing}&て{\middel}ば \mbox{(te{\middel}ba)}\\\hline
羽二重&feather + double = \textit{fine Japanese silk}&は{\middel}ぶた{\middel}え \mbox{(ha{\middel}buta{\middel}e)}\\\hline
三羽&three + feather = \textit{three (birds)}&サン{\middel}ば \mbox{(SAN{\middel}ba)}\\\hline
羽車&feather + car = \textit{portable shrine}&は{\middel}ぐるま \mbox{(ha{\middel}guruma)}\\\hline
\end{somme}

\begin{decomp}{8}
羽&は&鳥&に&固有&の&物&だ。\\\hline
はね&は&とり&に&コユウ&の&もの&だ\\\hline
feather&は&bird&に&inherent&の&thing&be/is\\\hline
\multicolumn{8}{|r|}{\textit{Feathers are peculiar to birds.}}\\\hline
\end{decomp}
\end{multicols}
\pagebreak
%%--------------------------------------------------------------------
%5--------------------------------------------------------------------

\begin{multicols}{2}
\begin{glyph}{Warm (温)}{warm}\label{glyph:warm}
Warm.
\end{glyph}
\gindex{Warm}{温}\strindex{12}{温}

\begin{comps}
\comprtwocone & 氵 + 日 + 皿 \\\hline
\comprthreecone & Water + Sun + Plate \\\hline
\end{comps}

\begin{remglyph}
\hooglig{Warm} \remcom{water} \remcom{today} on a \remcom{plate}.\\\hline
\end{remglyph}

\begin{prons}
\pronsrtwocone&\textbf{オン (ON)}&\textbf{あたた (atata)}\\\hline
\pronsrthreecone&---&---\\\hline
\end{prons}

\begin{remprons}
\hooglig{Warming} up \comon{online} is \comkun{a Tata} car.\\\hline
\end{remprons}

%{\middel}
% &  \mbox{()}\\\hline
\begin{somme}
温かい&\textit{warm}&あたた{\middel}かい \mbox{(atata{\middel}kai)}\\\hline
気温&spirit + warm = \textit{(air) temperature}&キ{\middel}オン \mbox{(KI{\middel}ON)}\\\hline
温泉&warm + springs = \textit{hot springs; spa}&オン{\middel}セン \mbox{(ON{\middel}SEN)}\\\hline
温度&warm + degree = \textit{temperature}&オン{\middel}ド \mbox{(ON{\middel}DO)}\\\hline
温室&warm + room = \textit{greenhouse}&オン{\middel}シツ \mbox{(ON{\middel}SHITSU)}\\\hline
室温&room + warm = \textit{room temperature}&シツ{\middel}オン \mbox{(SHITSU{\middel}ON)}\\\hline
温和&warm + harmony = \textit{clement; pleasant}&オン{\middel}ワ \mbox{(ON{\middel}WA)}\\\hline
体温&body + warm = \textit{body temperature}&タイ{\middel}オン \mbox{(TAI{\middel}ON)}\\\hline
\notyet{天然温泉}&nature + hot spring = \textit{natural hot spring}&テン{\middel}ネン{\middel}オン{\middel}セン \mbox{(TEN{\middel}NEN{\middel}ON{\middel}SEN)}\\\hline
\end{somme}
\end{multicols}

\begin{decomp}{5}
気温&が&急&に&下がった。\\\hline
キオン&が&キュウ&に&さがった\\\hline
temperature&が&hurry&に&to come down\\\hline
\multicolumn{5}{|r|}{\textit{The temperature has suddenly dropped.}}\\\hline
\end{decomp}
%%--------------------------------------------------------------------
%6--------------------------------------------------------------------
\begin{multicols}{2}
\begin{glyph}{Cold (寒)}{cold}\label{glyph:cold}
Cold.
\end{glyph}
\gindex{Cold}{寒}\strindex{12}{寒}

\begin{comps}
\comprtwocone&宀 + 井 + 一 + 八 + 氷\\\hline
\comprthreecone&Roof + Well + One + Eight + Ice\\\hline
\end{comps}

\begin{remglyph}
\hooglig{The cold} \remcom{roof} from \remcom{Town} \remcom{``1''} was thrown into the \remcom{volcano} of \remcom{ice}.\\\hline
\end{remglyph}

\begin{prons}
\pronsrtwocone&\textbf{カン (KAN)}&\textbf{さむ (samu)}\\\hline
\pronsrthreecone&---&---\\\hline
\end{prons}

\begin{remprons}
\hooglig{Cold} \comkun{Samurai} are very \comon{cunning}.\\\hline
\end{remprons}

%{\middel}
% &  \mbox{()}\\\hline
\begin{somme}
寒い&\textit{cold}&さむ{\middel}い \mbox{(samu{\middel}i)}\\\hline
寒気&cold + spirit = \textit{coldness; the cold}&カン{\middel}キ \mbox{(KAN{\middel}KI)}\\\hline
寒気&cold + spirit = \textit{chill; the shivers}&さむ{\middel}ケ \mbox{(samu{\middel}KE)}\\\hline
寒村&cold + village = \textit{deserted/poor village}&カン{\middel}ソン \mbox{(KAN{\middel}SON)}\\\hline
悪寒&evil + cold = \textit{shakes; ague}&オ{\middel}カン \mbox{(O{\middel}KAN)}\\\hline
寒冷前線&coldness + weather front = \textit{cold front}&カン{\middel}レイ{\middel}ゼン{\middel}セン \mbox{(KAN{\middel}REI{\middel}ZEN{\middel}SEN)}\\\hline
寒々&cold $(\times 2)$ = \textit{bleak; desolate; empty}&さむ{\middel}ざむ \mbox{(samu{\middel}zamu)}\\\hline
寒冷地&coldness + ground = \textit{cold region}&カン{\middel}レイ{\middel}チ \mbox{(KAN{\middel}REI{\middel}CHI)}\\\hline
寒心&cold + heart = \textit{deplorable; alarming}&カン{\middel}シン \mbox{(KAN{\middel}SHIN)}\\\hline
寒水石&cold water + stone = \textit{whitish marble}&カン{\middel}スイ{\middel}セキ \mbox{(KAN{\middel}SUI{\middel}SEKI)}\\\hline
\end{somme}

\begin{decomp}{4}
寒い&冬&が&来た。\\\hline
さむい&ふゆ&が&きた\\\hline
cold&winter&が&to come\\\hline
\multicolumn{4}{|r|}{\textit{Cold winter came on.}}\\\hline
\end{decomp}

\end{multicols}

\pagebreak
%%--------------------------------------------------------------------
%7--------------------------------------------------------------------
\begin{multicols}{2}
\begin{glyph}{Diligence (勉)}{diligence}\label{glyph:diligence}
Diligence.
\end{glyph}
\gindex{Diligence}{勉}\strindex{10}{勉}

\begin{comps}
\comprtwocone&免 + 力\\\hline
\comprthreecone&Exemption/Excuse/Dismissal + Strength\\\hline
\end{comps}

\begin{remglyph}
\hooglig{Diligence} will \remcom{exempt} you from mediocrity and give you \remcom{strength}.\\\hline
\end{remglyph}

\begin{prons}
\pronsrtwocone&\textbf{ベン (BEN)}&---\\\hline
\pronsrthreecone&---&---\\\hline
\end{prons}

\begin{remprons}
\hooglig{Diligence} will \comon{bend} the impossible.\\\hline
\end{remprons}

%{\middel}
% &  \mbox{()}\\\hline
\begin{somme}
勉学&diligence + study = \textit{diligent study}&ベン{\middel}ガク \mbox{(BEN{\middel}GAKU)}\\\hline
勉強&diligence + strong = \textit{studying}&ベン{\middel}キョウ \mbox{(BEN{\middel}KY\={O})}\\\hline
不勉強&not + studying = \textit{idleness}&フ{\middel}ベン{\middel}キョウ \mbox{(FU{\middel}BEN{\middel}KY\={O})}\\\hline
勉強部屋&studying + room = \textit{study room}&ベン{\middel}キョウ{\middel}べ{\middel}や \mbox{(BEN{\middel}KY\={O}{\middel}be{\middel}ya)}\\\hline
受験勉強&examination + studying = \textit{study for a test}&ジュ{\middel}ケン{\middel}ベン{\middel}キョウ \mbox{(JU{\middel}KEN{\middel}BEN{\middel}KY\={O})}\\\hline
社会勉強&society + studying = \textit{learning about the world}&シャ{\middel}カイ{\middel}ベン{\middel}キョウ \mbox{(SHA{\middel}KAI{\middel}BEN{\middel}KY\={O})}\\\hline
勉強会&studying + meeting = \textit{study group}&ベン{\middel}キョウ{\middel}カイ \mbox{(BEN{\middel}KY\={O}{\middel}KAI)}\\\hline
\notyet{勉強熱心}&studying + zeal = \textit{hardworking}&ベン{\middel}キョウ{\middel}ネッ{\middel}シン (BEN{\middel}KY\={O}{\middel}NES{\middel}SHIN)\\\hline
\end{somme}
\end{multicols}

\begin{decomp}{5}
彼女&は&勉学&に&熱中している。\\\hline
かのジョ&は&ベンガク&に&ネッチュウしている\\\hline
she&は&diligent study&に&zeal\\\hline
\multicolumn{5}{|r|}{\textit{She is intense in her study.}}\\\hline
\end{decomp}
%%--------------------------------------------------------------------
%8--------------------------------------------------------------------
\begin{multicols}{2}
\begin{glyph}{Technique (術)}{technique}\label{glyph:technique}
Technique.  Skill.
\end{glyph}
\gindex{Technique}{術}\strindex{11}{術}

\begin{comps}
\comprtwocone & 行 + (朮 = 木 + 丶) \\\hline
\comprthreecone & To go + (Dotted Ho = Tree + Dot) \\\hline
\end{comps}

\begin{remglyph}
\hooglig{Techniques} for \remcom{going} up \remcom{trees} to collect their \remcom{leaves}.\\\hline
\hooglig{Technique} \remcom{to go} with a \remcom{dotted ho}.\\\hline
\end{remglyph}

\begin{prons}
\pronsrtwocone&\textbf{ジュツ (JUTSU)}&---\\\hline
\pronsrthreecone&---&---\\\hline
\end{prons}

\begin{remprons}
\hooglig{Technique} is the English word for \comon{Jutsu}.\\\hline
\end{remprons}

%{\middel}
% &  \mbox{()}\\\hline
\begin{somme}
美術&beautiful + technique = \textit{art; the fine arts}&ビ{\middel}ジュツ \mbox{(BI{\middel}JUTSU)}\\\hline
手術&hand + technique = \textit{surgical operation}&シュ{\middel}ジュツ \mbox{(SHU{\middel}JUTSU)}\\\hline
術後&technique + after = \textit{after surgery}&ジュツ{\middel}ゴ \mbox{(JUTSU{\middel}GO)}\\\hline
話術&speak + technique = \textit{storytelling}&ワ{\middel}ジュツ \mbox{(WA{\middel}JUTSU)}\\\hline
学術&study + technique = \textit{academic pursuits}&ガク{\middel}ジュツ \mbox{(GAKU{\middel}JUTSU)}\\\hline
馬術&horse + technique = \textit{equestrian art}&バ{\middel}ジュツ \mbox{(BA{\middel}JUTSU)}\\\hline
占星術&occupy + star + technique = \textit{astrology}&セン{\middel}セイ{\middel}ジュツ \mbox{(SEN{\middel}SEI{\middel}JUTSU)}\\\hline
読心術&read + heart + technique = \textit{mind reading}&ドク{\middel}シン{\middel}ジュツ \mbox{(DOKU{\middel}SHIN{\middel}JUTSU)}\\\hline
\end{somme}

\begin{decomp}{5}
手術&が&必要&です&か。\\\hline
シュジュツ&が&ヒツヨウ&です&か\\\hline
operation&が&necessary&be/is&か\\\hline
\multicolumn{5}{|r|}{\textit{Do I need an operation?}}\\\hline
\end{decomp}
\end{multicols}

\pagebreak
%%--------------------------------------------------------------------
%9--------------------------------------------------------------------

\begin{multicols}{2}
\begin{glyph}{Personal (身)}{personal}\label{glyph:personal}
This 漢字 has to do with things ``personal'', in the sense of someone's \textit{own body} or a specific thing's \textit{characteristic}.
\end{glyph}
\gindex{Personal}{身}\strindex{7}{身}

\begin{comps}
\comprtwocone&身\\\hline
\comprthreecone&Body\\\hline
\end{comps}

\begin{remglyph}
Part of the 漢字 radicals (\#{158}).\\\hline
\end{remglyph}

\begin{prons}
\pronsrtwocone&\textbf{シン (SHIN)}&\textbf{み (mi)}\\\hline
\rye{3}{み predominates in the first position, and シン in the second.}
\pronsrthreecone&---&---\\\hline
\end{prons}

\begin{remprons}
\hooglig{The person} in the \remcom{mirror} has a nice \remcom{sheen}.\\\hline
\end{remprons}

%{\middel}
% &  \mbox{()}\\\hline
\begin{somme}
身体&personal + body = \textit{the body; physical}&シン{\middel}タイ \mbox{(SHIN{\middel}TAI)}\\\hline
出身&exit + personal = \textit{one's home-town}&シュッ{\middel}シン \mbox{(SHUS{\middel}SHIN)}\\\hline
中身&middle + personal = \textit{contents}&なか{\middel}み \mbox{(naka{\middel}mi)}\\\hline
自身&self + personal = \textit{one's self; oneself}&ジ{\middel}シン \mbox{(JI{\middel}SHIN)}\\\hline
身長&personal + long = \textit{body height; stature}&シン{\middel}チョウ \mbox{(SHIN{\middel}CH\={O})}\\\hline
全身&complete + personal = \textit{whole body; full-length}&ゼン{\middel}シン \mbox{(ZEN{\middel}SHIN)}\\\hline
身分&personal + part = \textit{social status}&み{\middel}ブン \mbox{(mi{\middel}BUN)}\\\hline
心身&heart + personal = \textit{mind and body}&シン{\middel}シン \mbox{(SHIN{\middel}SHIN)}\\\hline
身近&personal + near = \textit{familiar}&み{\middel}ぢか \mbox{(mi{\middel}jika)}\\\hline
受身&receive + personal = \textit{the defensive; passiveness}&うけ{\middel}み \mbox{(uke{\middel}mi)}\\\hline
\end{somme}

\begin{decomp}{5}
京都&の&出身&です&か。\\\hline
キョウト&の&シュッシン&です&か\\\hline
Ky\={o}to&の&hometown&be/is&か\\\hline
\multicolumn{5}{|r|}{\textit{Are you from Ky\={o}to?}}\\\hline
\end{decomp}

\end{multicols}
%%--------------------------------------------------------------------
%10-------------------------------------------------------------------
\begin{multicols}{2}
\begin{glyph}{Beginning (初)}{beginning}\label{glyph:beginning}
Beginning.
\end{glyph}
\gindex{Beginning}{初}\strindex{7}{初}

\begin{comps}
\comprtwocone&衤 + 刀\\\hline
\comprthreecone&Clothing + Sword\\\hline
\end{comps}

\begin{remglyph}
\hooglig{Begin} \remcom{clothing} the \remcom{sword}.\\\hline
\end{remglyph}

\begin{prons}
\pronsrtwocone&\textbf{ショ (SHO)}&\textbf{はじ、はつ (haji, hatsu)}\\\hline
\rye{3}{ショ is the most common of these three.  はじ appears only in \mbox{初めて}.  はつ has a more poetic meaning.}
\pronsrthreecone&---&そ (so)\\\hline
\end{prons}

\begin{remprons}
\hooglig{Beginning} of \comkun{Hajj} we \comon{shop} around for \ucomkun{soft} \comkun{huts} to stay in.\\\hline
\end{remprons}

%{\middel}
% &  \mbox{()}\\\hline
\begin{somme}
初め&\textit{beginning; outset}&はじ{\middel}め \mbox{(haji{\middel}me)}\\\hline
初日&beginning + sun = \textit{first/opening day}&ショ{\middel}ニチ \mbox{(SHO{\middel}NICHI)}\\\hline
初耳&beginning + ear = \textit{hearing something for the first time}&はつ{\middel}みみ \mbox{(hatsu{\middel}mimi)}\\\hline
初代&beginning + substitute = \textit{$1^{st}$ generation; founder}&ショ{\middel}ダイ \mbox{(SHO{\middel}DAI)}\\\hline
初回&beginning + rotate = \textit{initial attempt}&ショ{\middel}カイ \mbox{(SHO{\middel}KAI)}\\\hline
初心者&naive + individual = \textit{beginner}&ショ{\middel}シン{\middel}シャ \mbox{(SHO{\middel}SHIN{\middel}SHA)}\\\hline
初歩&beginning + walk = \textit{basics}&ショ{\middel}ホ \mbox{(SHO{\middel}HO)}\\\hline
初顔合わせ&first + face + meet = \textit{first meeting}&はつ{\middel}かお{\middel}あ{\middel}わせ \mbox{(hatsu{\middel}kao{\middel}a{\middel}wase)}\\\hline
\end{somme}

\begin{decomp}{4}
間もなく&初雪&が&降った。\\\hline
まもなく&はつゆき&が&ふった\\\hline
soon&first snowfall&が&to descend\\\hline
\multicolumn{4}{|r|}{\textit{The first snow came before long.}}\\\hline
\end{decomp}
\end{multicols}

\pagebreak
%%--------------------------------------------------------------------
%11-------------------------------------------------------------------

\begin{multicols}{2}
\begin{glyph}{Distress (困)}{distress}\label{glyph:distress}
Distress.
\end{glyph}
\gindex{Distress}{困}\strindex{7}{困}

\begin{comps}
\comprtwocone&囗 + 木 \\\hline
\comprthreecone&Enclosure/Prison + Tree\\\hline
\end{comps}

\begin{remglyph}
\hooglig{Distress} in the \remcom{prison} of \remcom{trees}.\\\hline
\end{remglyph}

\begin{prons}
\pronsrtwocone&\textbf{コン (KON)}&\textbf{こま (koma)}\\\hline
\pronsrthreecone&---&---\\\hline
\end{prons}

\begin{remprons}
\hooglig{Distress} from all the \comon{cons} resulted in a \comkun{coma}.\\\hline
\end{remprons}

%{\middel}
% &  \mbox{()}\\\hline
\begin{somme}
困る&\textit{be distressed; suffer}&こま{\middel}る \mbox{(koma{\middel}ru)}\\\hline
困難&distress + difficult = \textit{difficulty}&コン{\middel}ナン \mbox{(KON{\middel}NAN)}\\\hline
困り物&distress + thing = \textit{good-for-nothing; nuisance}&こま{\middel}り{\middel}もの \mbox{(koma{\middel}ri{\middel}mono)}\\\hline
生活に困る&livelihood + distress = \textit{living in want}&せい{\middel}かつ{\middel}に{\middel}こま{\middel}る (SEI{\middel}KATSU{\middel}ni{\middel}koma{\middel}ru)\\\hline
入手困難&acquisition + difficulty = \textit{difficult to get}&ニュウ{\middel}シュ{\middel}コン{\middel}ナン (NY\={U}{\middel}SHU{\middel}KON{\middel}NAN)\\\hline
困った人&unmanageable + person = \textit{difficult person}&こまっ{\middel}たひと \mbox{(komat{\middel}ta{\middel}hito)}\\\hline
困り顔&distress + face = \textit{trouble expression}&かま{\middel}り{\middel}かお \mbox{(koma{\middel}ri{\middel}kao)}\\\hline
困り事&distress + matter = \textit{problem; trouble}&こま{\middel}り{\middel}ごと \mbox{(koma{\middel}ri{\middel}goto)}\\\hline
\end{somme}

\begin{decomp}{4}
新しい&困難&が&生じた。\\\hline
あたらしい&コンナン&が&ショウじた\\\hline
new&difficulty&が&to produce\\\hline
\multicolumn{4}{|r|}{\textit{A new difficulty has arisen.}}\\\hline
\end{decomp}

\end{multicols}
%%--------------------------------------------------------------------
%12-------------------------------------------------------------------
\begin{multicols}{2}
\begin{glyph}{Put (置)}{put}\label{glyph:put}
Put.  Place.  Leave behind.\\\hline
Ensure you distinguish between this 漢字 and 直 (\textit{Straight}, Entry \ref{glyph:entry}, p. \pageref{glyph:entry}), 真 (\textit{True}, Entry \ref{glyph:true}, p. \pageref{glyph:true}).
\end{glyph}
\gindex{Put}{置}\strindex{13}{置}

\begin{comps}
\comprtwocone&直 + 罒\\\hline
\comprthreecone&Straight (p. \pageref{glyph:straight}) + Net\\\hline
\end{comps}

\begin{remglyph}
\hooglig{Put} the \remcom{net} \remcom{straight}.\\\hline
\end{remglyph}

\begin{prons}
\pronsrtwocone&\textbf{チ (CHI)}&\textbf{お (o)}\\\hline
\rye{3}{お is the reading in the first position, with チ more common in the second.}
\pronsrthreecone&---&---\\\hline
\end{prons}

\begin{remprons}
\hooglig{Put} behind \comon{cheap} \comkun{obligations}.\\\hline
\end{remprons}

%{\middel}
% &  \mbox{()}\\\hline
\begin{somme}
\rye{3}{This 漢字 can absorb its accompanying ひらがな.}
置く&\textit{put; place; leave}&お{\middel}く \mbox{(o{\middel}ku)}\\\hline
置場&put + place = \textit{place (for something)}&おき{\middel}ば \mbox{(oki{\middel}ba)}\\\hline
配置&distribute + put = \textit{arrangement; layout}&ハイ{\middel}チ \mbox{(HAI{\middel}CHI)}\\\hline
物置&thing + put = \textit{storeroom}&もの{\middel}おき \mbox{(mono{\middel}oki)}\\\hline
前置き&before + put = \textit{preamble; intro}&まえ{\middel}お{\middel}き \mbox{(mae{\middel}o{\middel}ki)}\\\hline
置き去り&put + depart = \textit{abandonment}&お{\middel}き{\middel}ざ{\middel}り \mbox{(o{\middel}ki{\middel}za{\middel}ri)}\\\hline
安置&ease + put = \textit{enshrinement}&アン{\middel}チ \mbox{(AN{\middel}CHI)}\\\hline
買い置き&buy + put = \textit{stock}&か{\middel}い{\middel}お{\middel}き \mbox{(ka{\middel}i{\middel}o{\middel}ki)}\\\hline
定置&determine + put = \textit{fixed}&テイ{\middel}チ \mbox{(TEI{\middel}CHI)}\\\hline
\end{somme}
\end{multicols}

\begin{decomp}{7}
彼女&の&息子&は&西ドイツ&に&配置されている。\\\hline
かのジョ&の&むすこ&は&にちドイツ&に&ハイチされている\\\hline
she&の&son&は&West Germany&に&deployment\\\hline
\multicolumn{7}{|r|}{\textit{Her son is stationed in West Germany.}}\\\hline
\end{decomp}

\pagebreak
%%--------------------------------------------------------------------
%13-------------------------------------------------------------------
\begin{multicols}{2}
\begin{glyph}{Solo (個)}{solo}\label{glyph:solo}
This 漢字 emphasizes the separateness of individual objects (including people).\\\hline
It is also used as a general counter for ``pieces''.
\end{glyph}
\gindex{Solo}{個}\strindex{10}{個}

\begin{comps}
\comprtwocone&亻 + 固\\\hline
\comprthreecone&Man/Person + Hard (p. \pageref{glyph:hard})\\\hline
\end{comps}

\begin{remglyph}
\hooglig{Solo} \remcom{people} have it \remcom{hard}.\\\hline
\end{remglyph}

\begin{prons}
\pronsrtwocone&\textbf{コ (KO)}&---\\\hline
\pronsrthreecone&---&---\\\hline
\end{prons}

\begin{remprons}
\hooglig{Solo} \comon{Cop}.\\\hline
\end{remprons}

%{\middel}
% &  \mbox{()}\\\hline
\begin{somme}
個別&solo + different = \textit{particular case}&コ{\middel}ベツ \mbox{(KO{\middel}BETSU)}\\\hline
別個&different + solo = \textit{another; different}&ベツ{\middel}コ \mbox{(BETSU{\middel}KO)}\\\hline
一個&one + solo = \textit{piece; fragment}&イッ{\middel}コ \mbox{(IK{\middel}KO)}\\\hline
{\diff}個所&solo + location = \textit{spot}&カ{\middel}ショ \mbox{(KA{\middel}SHO)}\\\hline
個体&solo + body = \textit{specimen}&コ{\middel}タイ \mbox{(KO{\middel}TAI)}\\\hline
個人的意見&personal + view = \textit{personal opinion}&コ{\middel}ジン{\middel}テキ{\middel}イ{\middel}ケン (KO{\middel}JIN{\middel}TEKI{\middel}I{\middel}KEN)\\\hline
何個&what + solo = \textit{how many pieces?}&なん{\middel}コ \mbox{(nan{\middel}KO)}\\\hline
個人事業&individual + business = \textit{personal business}&コ{\middel}ジン{\middel}ジ{\middel}ギョウ \mbox{(KO{\middel}JIN{\middel}JI{\middel}GY\={O})}\\\hline
個人特定&individual + specific = \textit{pinpointing so/sth}&コ{\middel}ジン{\middel}トク{\middel}テイ \mbox{(KO{\middel}JIN{\middel}TOKU{\middel}TEI)}\\\hline
\end{somme}
\end{multicols}

\begin{decomp}{6}
個人的&な&話&に&受け取らない&で。\\\hline
コジンテキ&な&はなし&に&うけとらない&で\\\hline
personal&な&argument&に&to not receive&で\\\hline
\multicolumn{6}{|r|}{\textit{Don't take it personally.}}\\\hline
\end{decomp}
%%--------------------------------------------------------------------
%14-------------------------------------------------------------------
\begin{multicols}{2}
\begin{glyph}{Become (成)}{become}\label{glyph:become}
Become.  Turn into.  Reach.\\\hline
Pay close attention to the stroke order of this 漢字.
\end{glyph}
\gindex{Become}{成}\strindex{6}{成}

\begin{comps}
\comprtwocone&勹 +　戈\\\hline
\comprthreecone&Wrap/Embrace + Spear/Halberd\\\hline
\end{comps}

\begin{remglyph}
\hooglig{Become} a \remcom{foreskin} mutilator with a \remcom{spear}.\\\hline
\end{remglyph}

\begin{prons}
\pronsrtwocone&\textbf{セイ (SEI)}&\textbf{な (na)}\\\hline
\rye{3}{The 訓読み can incorporate its ひらがな ending in certain compounds.}
\pronsrthreecone&ジョウ (J\={O})&---\\\hline
\end{prons}

\begin{remprons}
\hooglig{Become} \comkun{numb} in your \ucomon{jaw} from all the \comon{sailing}.\\\hline
\end{remprons}

%{\middel}
% &  \mbox{()}\\\hline
\begin{somme}
成る&\textit{become}&な{\middel}る \mbox{(na{\middel}ru)}\\\hline
成年&become + year = \textit{adult; majority}&セイ{\middel}ネン \mbox{(SEI{\middel}NEN)}\\\hline
成立&become + stand = \textit{formation}&セイ{\middel}リツ \mbox{(SEI{\middel}RITSU)}\\\hline
成長&become + long = \textit{growth; development}&セイ{\middel}チョウ \mbox{(SEI{\middel}CHOU)}\\\hline
行き成&go + become = \textit{suddenly; without warning}&い{\middel}き{\middel}なり \mbox{(i{\middel}ki{\middel}nari)}\\\hline
結成&bind + become = \textit{formation; combination}&ケッ{\middel}セイ \mbox{(KES{\middel}SEI)}\\\hline
合成&match + become = \textit{composition; synthesis}&ゴウ{\middel}セイ \mbox{(G\={O}{\middel}SEI)}\\\hline
達成&achieve + become = \textit{accomplishment; realisation}&タッ{\middel}セイ \mbox{(TAS{\middel}SEI)}\\\hline
\end{somme}

\begin{decomp}{3}
早く&良く&成れ。\\\hline
はやく&よく&なれ\\\hline
soon&well&to become\\\hline
\multicolumn{3}{|r|}{\textit{I hope you will get well soon.}}\\\hline
\end{decomp}

\end{multicols}
\pagebreak
%%--------------------------------------------------------------------
%15-------------------------------------------------------------------

\begin{multicols}{2}
\begin{glyph}{Work (働)}{work}\label{glyph:work}
Work.
\end{glyph}
\gindex{Work}{働}\strindex{13}{働}

\begin{comps}
\comprtwocone&亻 + 動\\\hline
\comprthreecone&Man/Person + Move (p. \pageref{glyph:move})\\\hline
\end{comps}

\begin{remglyph}
\hooglig{Working} \remcom{people} \remcom{move}.\\\hline
\end{remglyph}

\begin{prons}
\pronsrtwocone&---&\textbf{はたら (hatara)}\\\hline
\pronsrthreecone&ドウ (D\={O})&---\\\hline
\end{prons}

\begin{remprons}
\hooglig{Work} hard delivering oxygen from \ucomon{door} to door in ``\comkun{Hataraku} Saibou''.\\\hline
{\warn}\textit{Hataraku Saibou} is an アニメ series about the human body.\\\hline
\end{remprons}

%{\middel}
% &  \mbox{()}\\\hline
\begin{somme}
働く&\textit{work}&はたら{\middel}く \mbox{(hatara{\middel}ku)}\\\hline
働き口&work + mouth = \textit{job; position}&はたら{\middel}き{\middel}ぐち \mbox{(hatara{\middel}ki{\middel}guchi)}\\\hline
\end{somme}

\begin{decomp}{3}
何時&働いてる&の？\\\hline
いつ&はたらいてる&の\\\hline
when&to work&の\\\hline
\multicolumn{3}{|r|}{\textit{When do you work?}}\\\hline
\end{decomp}

\end{multicols}

%%--------------------------------------------------------------------
%16-------------------------------------------------------------------
\begin{multicols}{2}
\begin{glyph}{Hall (堂)}{hall}\label{glyph:hall}
This 漢字 refers to a \textit{larger-sized hall}, usually in a public building.
\end{glyph}
\gindex{Hall}{堂}\strindex{11}{堂}

\begin{comps}
\comprtwocone&⺌小 + 冖 + 口 + 土\\\hline
\comprthreecone&Small + Cover/Crown + Mouth + Earth\\\hline
\rye{2}{Classroom = ⺌小 + 冖 = Small + Cover/Crown.  \hooglig{Classrooms} are so \remcom{small} you can \remcom{cover} them.}
\end{comps}

\begin{remglyph}
\hooglig{Halls} that are \remcom{classrooms} for \remcom{vampires} are \remcom{dirty}.\\\hline
\end{remglyph}

\begin{prons}
\pronsrtwocone&\textbf{ドウ (D\={O})}&---\\\hline
\pronsrthreecone&---&---\\\hline
\end{prons}

\begin{remprons}
\hooglig{Hallway's} \comon{door}.\\\hline
\end{remprons}

%{\middel}
% &  \mbox{()}\\\hline
\begin{somme}
食堂&eat + hall = \textit{dining hall; cafeteria}&ショク{\middel}ドウ \mbox{(SHOKU{\middel}D\={O})}\\\hline
本堂&main + hall = \textit{main hall; main temple}&ホン{\middel}ドウ \mbox{(HON{\middel}D\={O})}\\\hline
\end{somme}
\end{multicols}

\begin{decomp}{6}
食堂&は&何時&に&開きます&か。\\\hline
ショクドウ&は&なんジ&に&あきます&か\\\hline
dining room&は&what time&に&to open&か\\\hline
\multicolumn{6}{|r|}{\textit{What time does the dining room open?}}\\\hline
\end{decomp}

\pagebreak
%%--------------------------------------------------------------------
%17-------------------------------------------------------------------

\begin{multicols}{2}
\begin{glyph}{Association (組)}{association}\label{glyph:association}
Association.  Set.\\\hline
This 漢字 has to do with \textit{things coming together} or \textit{being put together} in groups.
\end{glyph}
\gindex{Association}{組}\strindex{11}{組}

\begin{comps}
\comprtwocone &糸 + 且\\\hline
\comprthreecone &Thread + Besides/Moreover (\ref{glyph:besides})\\\hline
\end{comps}

\begin{remglyph}
\hooglig{Set} the \remcom{thread} \remcom{beside} the needle.\\\hline
\end{remglyph}

\begin{prons}
\pronsrtwocone & --- & \textbf{く、くみ (ku, kumi)}\\\hline
\rye{3}{The 訓読み become voiced in the second position according to the rules we learned for 付 in Entry \ref{glyph:attach}.}
\pronsrthreecone & ソ (SO) & --- \\\hline
\end{prons}

\begin{remprons}
\hooglig{Association} \ucomon{solved} the \comkun{cumin} \comkun{coup}.\\\hline
{\warn}\textit{Cumin} is a spice made from dried seeds.\\\hline
\end{remprons}

%{\middel}
% &  \mbox{()}\\\hline
\begin{somme}
\rye{3}{Note the mix of 音読み and 訓読み in the final compound.}
組む & \textit{put together; form} & く{\middel}む \mbox{(ku{\middel}mu)}\\\hline
組合 & association + match = \textit{association; (labor) union} & くみ{\middel}あい \mbox{(kumi{\middel}ai)}\\\hline
番組 & order + association = \textit{a program} & バン{\middel}ぐみ \mbox{(BAN{\middel}gumi)}\\\hline
\end{somme}

\begin{decomp}{5}
番組&は&国歌&で&終了する。\\\hline
バンぐみ&は&コッカ&で&シュウリョウする\\\hline
program&は&national anthem&で&end\\\hline
\multicolumn{5}{|r|}{\textit{The program will finish with the national anthem.}}\\\hline
\end{decomp}

\end{multicols}

%%--------------------------------------------------------------------
%18-------------------------------------------------------------------
\begin{multicols}{2}
\begin{glyph}{Arise (起)}{arise}\label{glyph:arise}
Arise.  Get up.  Happen.\\\hline
The 漢字 for ``run'' always stretches out like this when appearing on the left of another character as a component.
\end{glyph}
\gindex{Arise}{起}\strindex{10}{起}

\begin{comps}
\comprtwocone & 走 + 己 \\\hline
\comprthreecone & Run + Oneself \\\hline
\end{comps}

\begin{remglyph}
\hooglig{Arise} \remcom{oneself} and start \remcom{running}.\\\hline
\end{remglyph}

\begin{prons}
\pronsrtwocone & \textbf{キ (KI)} & \textbf{お (o)}\\\hline
\pronsrthreecone & --- & --- \\\hline
\end{prons}

\begin{remprons}
\hooglig{Arise} for the \comon{kill} you are \comkun{obligated} to carry out.\\\hline
\end{remprons}

%{\middel}
% &  \mbox{()}\\\hline
\begin{somme}
\rye{3}{The first two examples below function as a normal intransitive/transitive verb pair.}
起こる (intr.) & \textit{to happen; come to pass} & お{\middel}こる \mbox{(o{\middel}koru)}\\\hline
起こす (tr.) & \textit{to bring out; raise up (sth.)} & お{\middel}こす \mbox{(o{\middel}kosu)}\\\hline
おきる (intr.) & \textit{to get up; rise} & お{\middel}きる \mbox{(o{\middel}kiru)}\\\hline
起点 & arise + point = \textit{starting point} & キ{\middel}テン \mbox{(KI{\middel}TEN)}\\\hline
\end{somme}

\begin{decomp}{3}
起きる&時間&よ。\\\hline
おきる&ジカン&よ\\\hline
to get up&time&よ\\\hline
\multicolumn{3}{|r|}{\textit{Time to get up.}}\\\hline
\end{decomp}

\end{multicols}

\pagebreak
%%--------------------------------------------------------------------
%19-------------------------------------------------------------------
\begin{multicols}{2}
\begin{glyph}{Dispute (争)}{dispute}\label{glyph:dispute}
Dispute.  Argue.
\end{glyph}
\gindex{Dispute}{争}\strindex{6}{争}

\begin{comps}
\comprtwocone & 勹 + 亅 + 彐彑 \\\hline
\comprthreecone & Wrap/Embrace + Hook + Pig's Head/Snout \\\hline
\end{comps}

\begin{remglyph}
\hooglig{Dispute} about who should \remcom{wrap} the \remcom{skewered} \remcom{pig's head}.\\\hline
\end{remglyph}

\begin{prons}
\pronsrtwocone & \textbf{ソウ (S\={O})} & \textbf{あらそ (araso)}\\\hline
\pronsrthreecone & --- & --- \\\hline
\end{prons}

\begin{remprons}
\hooglig{Dispute} \comkun{arouses} \comon{soreness} of heart.\\\hline
\end{remprons}

%{\middel}
% &  \mbox{()}\\\hline
\begin{somme}
争う & \textit{dispute; contend} & あらそ{\middel}う \mbox{(araso{\middel}u)}\\\hline
争点 & dispute + point = \textit{point of dispute} & ソウ{\middel}テン \mbox{(S\={O}{\middel}TEN)}\\\hline
内争 & inside + dispute = \textit{internal dispute} & ナイ{\middel}ソウ \mbox{(NAI{\middel}S\={O})}\\\hline
\end{somme}
\end{multicols}

\begin{decomp}{6}
あなた&と&言い争って&も&無駄&だ。\\\hline
あなた&と&いいあらそって&も&むだ&だ\\\hline
you&と&to argue&も&pointless&be / is\\\hline
\multicolumn{6}{|r|}{\textit{It is no good arguing with you.}}\\\hline
\end{decomp}


%%--------------------------------------------------------------------
%20-------------------------------------------------------------------
\begin{multicols}{2}
\begin{glyph}{Wind (風)}{wind}\label{glyph:wind}
Wind (as in breeze).\\\hline
An important secondary meaning relates to the ideas of \textit{style} and \textit{appearance}; the final compound provides an example.
\end{glyph}
\gindex{Wind}{風}\strindex{9}{風}

\begin{comps}
\comprtwocone & 風 \\\hline
\comprthreecone & Wind \\\hline
\end{comps}

\begin{remglyph}
Part of the 漢字 radicals (\#{182}).\\\hline
\end{remglyph}

\begin{prons}
\pronsrtwocone & \textbf{フウ (F\={U})} & \textbf{かざ、かぜ (kaza, kaze)}\\\hline
\rye{3}{フウ is the most common of the 3 readings, while かざ　appears only in the first position.  かぜ is the reading used when the 漢字 stand on its own.  かぜ rarely occurs in the first position, but it is as common as フウ when this character is the last in a compound.\\\hline
フウ　is the reading used in compounds relating to style and appearance.}
\pronsrthreecone & フ (FU) & --- \\\hline
\rye{3}{フ is used in \textit{bath}: 風呂\notcov (フ{\middel}ロ/FU{\middel}RO), from the 漢字 `wind' and `spine'.  (`Spine' is not covered in this book.)}
\end{prons}

\begin{remprons}
\hooglig{The wind} in \comkun{Kazakstan} made a \ucomon{foofy} of its \comkun{Kazekage}, which is some \comon{food} for thought.\\\hline
{\warn}\textit{Kazekage} is a fictional character in the ナルト アニメ.\\\hline
\end{remprons}

%{\middel}
% &  \mbox{()}\\\hline
\begin{somme}
風 & \textit{wind} & かぜ \mbox{(kaze)}\\\hline
風車 & wind + car = \textit{windmill} & かざ{\middel}ぐるま \mbox{(kaza{\middel}guruma)}\\\hline
台風 & platform + wind = \textit{typhoon} & タイ{\middel}フウ \mbox{(TAI{\middel}F\={U})}\\\hline
日本風 & Japan + wind (style) = \textit{Japanese Style} & ニ{\middel}ホン{\middel}フウ \mbox{(NI{\middel}HON{\middel}F\={U})}\\\hline
\end{somme}

\begin{decomp}{4}
昨日&台風&が&来ました。\\\hline
きのう&タイフウ&が&きました\\\hline
yesterday&typhoon&が&to come\\\hline
\multicolumn{4}{|r|}{\textit{We had a storm yesterday.}}\\\hline
\end{decomp}

\end{multicols}
\pagebreak
%%--------------------------------------------------------------------
%21-------------------------------------------------------------------

\begin{multicols}{2}
\begin{glyph}{Pain (痛)}{pain}\label{glyph:pain}
Pain.
\end{glyph}
\gindex{Pain}{痛}\strindex{12}{痛}

\begin{comps}
\comprtwocone & 疒 + マ + 用 \\\hline
\comprthreecone & Sickness + MA + Utilize \\\hline
\end{comps}

\begin{remglyph}
\hooglig{Pain} \remcom{utilizes} \remcom{sickness} to hurt \remcom{mother}.\\\hline
\end{remglyph}

\begin{prons}
\pronsrtwocone & \textbf{ツウ (TS\={U})} & \textbf{いた (ita)}\\\hline
\pronsrthreecone & --- & --- \\\hline
\end{prons}

\begin{remprons}
\hooglig{Painful} \comkun{Italy} and \comon{Tsuu'Tina} clashes.\\\hline
\end{remprons}

%{\middel}
% &  \mbox{()}\\\hline
\begin{somme}
\rye{3}{The final compound below uses a less common 音読み for 頭 (\ref{glyph:head}).}
痛い & \textit{painful} & いた{\middel}い \mbox{(ita{\middel}i)}\\\hline
痛み止め & pain + stop = \textit{painkiller} & いた{\middel}み　ど{\middel}め \mbox{(ita{\middel}mi do{\middel}me)}\\\hline
頭痛 & head + pain = \textit{headache} & ズ{\middel}ツウ \mbox{(ZU{\middel}TS\={U})}\\\hline
\end{somme}

\begin{decomp}{4}
少し&頭痛&が&します。\\\hline
すこし&ズツウ&が&します\\\hline
little&headache&が&to be\\\hline
\multicolumn{4}{|r|}{\textit{I have a slight headache.}}\\\hline
\end{decomp}

\end{multicols}

%%--------------------------------------------------------------------
%22-------------------------------------------------------------------
\begin{multicols}{2}
\begin{glyph}{Officer (吏)}{officer}\label{glyph:officer}
An officer, or official, typically one of lower standing.\\\hline
We will find more use for this 漢字 as a component than as a character on its own.
\end{glyph}
\gindex{Officer}{吏}\strindex{6}{吏}

\begin{comps}
\comprtwocone & 一 + 史 \\\hline
\comprthreecone & One + History (\ref{glyph:history}) \\\hline
\end{comps}

\begin{remglyph}
That \hooglig{Officer} is \remcom{one} going down in \remcom{history}.\\\hline
\end{remglyph}

\begin{prons}
\pronsrtwocone & \textbf{リ (RI)} & --- \\\hline
\pronsrthreecone & --- & --- \\\hline
\end{prons}

\begin{remprons}
\hooglig{Officers} \comon{leaping} for joy.\\\hline
\end{remprons}

%{\middel}
% &  \mbox{()}\\\hline
\begin{somme}
官吏 & official + officer = \textit{government officer/official} & カン{\middel}リ \mbox{(KAN{\middel}RI)}\\\hline
\end{somme}

\begin{decomp}{4}
邪&な&官吏&だ。\\\hline
よこしま&な&カンリ&だ\\\hline
wicked / evil&な&official&be / is\\\hline
\multicolumn{4}{|r|}{\textit{He is a corrupt official.}}\\\hline
\end{decomp}

\end{multicols}

\pagebreak
%%--------------------------------------------------------------------
%23-------------------------------------------------------------------

\begin{multicols}{2}
\begin{glyph}{Sleep (寝)}{sleep}\label{glyph:sleep}
Sleep.
\end{glyph}
\gindex{Sleep}{寝}\strindex{13}{寝}

\begin{comps}
\comprtwocone & 宀 + 冖 + 彐彑 + 又 + 爿 \\\hline
\comprthreecone & Roof + Cover + Pig Head/Snout + Again + Bed \\\hline
\end{comps}

\begin{remglyph}
\hooglig{Sleeping} \remcom{again} with my \remcom{piggish head} under the \remcom{cover} of the \remcom{roof} in my \remcom{bed}.\\\hline
\end{remglyph}

\begin{prons}
\pronsrtwocone & \textbf{シン (SHIN)} & \textbf{ね (ne)}\\\hline
\rye{3}{The 訓読み is the more common reading here.}
\pronsrthreecone & --- & --- \\\hline
\end{prons}

\begin{remprons}
\hooglig{Sleeping} \comkun{necromancers} deathly \comon{sheens} will put you immediately to sleep.\\\hline
\end{remprons}

%{\middel}
% &  \mbox{()}\\\hline
\begin{somme}
寝る & \textit{sleep} & ね{\middel}る \mbox{(ne{\middel}ru)}\\\hline
昼寝 & daytime + sleep = \textit{afternoon nap} & ひる{\middel}ね \mbox{(hiru{\middel}ne)}\\\hline
寝室 & sleep + room = \textit{bedroom} & シン{\middel}シツ \mbox{(SHIN{\middel}SHITSU)}\\\hline
\end{somme}

\begin{decomp}{5}
寝室&で&物音&が&聞こえた。\\\hline
シンシツ&で&ものおと&が&きこえた\\\hline
bedroom&で&sound&が&to be heard\\\hline
\multicolumn{5}{|r|}{\textit{I heard a noise in the bedroom.}}\\\hline
\end{decomp}

\end{multicols}

%%--------------------------------------------------------------------
%24-------------------------------------------------------------------
\begin{multicols}{2}
\begin{glyph}{Checkpoint (関)}{checkpoint}\label{glyph:checkpoint}
Checkpoint.  Barrier.
\end{glyph}
\gindex{Checkpoint}{関}\strindex{14}{関}

\begin{comps}
\comprtwocone & 門 + (关 = 艹 + 天)\\\hline
\comprthreecone & Gate + (Relation = Grass + Heaven (\ref{glyph:heaven}))\\\hline
\end{comps}

\begin{remglyph}
\hooglig{The checkpoint} of \remcom{heaven's} \remcom{gate} was marked with beautiful green \remcom{grass}.\\\hline
\hooglig{Checkpoint} of the \remcom{gate} I have a \remcom{relationship} with.\\\hline
\end{remglyph}

\begin{prons}
\pronsrtwocone & \textbf{カン (KAN)} & \textbf{せき (seki)}\\\hline
\pronsrthreecone & --- & --- \\\hline
\end{prons}

\begin{remprons}
\hooglig{Checkpoint} made in the \comon{cunning} game of \comkun{Sekiro}: Shadows die twice.\\\hline
\end{remprons}

%{\middel}
% &  \mbox{()}\\\hline
\begin{somme}
関 & \textit{checkpoint; barrier} & せき \mbox{(seki)}\\\hline
機関 & mechanism + checkpoint = \textit{engine; organization} & キ{\middel}カン \mbox{(KI{\middel}KAN)}\\\hline
関係 & checkpoint + connection = \textit{relationship; connection} & カン{\middel}ケイ \mbox{(KAN{\middel}KEI)}\\\hline
\end{somme}

\begin{decomp}{5}
誰&か&玄関&に&要る。\\\hline
だれ&か&ゲンカン&に&いる\\\hline
who&か&entranceway&に&to need / want\\\hline
\multicolumn{5}{|r|}{\textit{Someone is at the door.}}\\\hline
\end{decomp}
\end{multicols}
\pagebreak
%%--------------------------------------------------------------------